\documentclass[a4paper,14pt]{article}

\usepackage[T2A]{fontenc}
\usepackage[utf8]{inputenc}
\usepackage[russian]{babel}

\usepackage[left=30mm,right=15mm,top=20mm,bottom=20mm]{geometry}

\usepackage{setspace}
\onehalfspacing

\usepackage{graphicx}
\usepackage{float}
\usepackage{caption}
\usepackage{longtable}
\usepackage{booktabs}
\usepackage{multirow}
\usepackage{amsmath,amssymb}
\usepackage{hyperref}

\numberwithin{figure}{section}
\numberwithin{table}{section}

\begin{document}

\section{Использование механизма внимания в задачах сегментации для построений динамических зон интереса}

\subsection{Аннотация}

Задача семантической сегментации изображений традиционно решается свёрточными нейронными сетями (CNN), 
однако ограниченное рецептивное поле свёрток не позволяет им улавливать глобальный контекст сцены, 
и сегментация фактически сводится к классификации отдельных пикселей по их локальному окружению. 
Сегментаторы на основе визуальных трансформеров (ViT) устраняют этот недостаток за счёт механизма самовнимания, 
связывающего каждый фрагмент изображения со всеми остальными, но формируют лишь одномасштабные признаки. 
Архитектура SegFormer строит иерархию многоуровневых признаков и с помощью механизма эффективного самовнимания (Efficient Self-Attention)
 снижает вычислительную сложность с $O(N^2)$ до $O(N^2/R)$, где $R$ - коэффициент сокращения последовательности. 
 Тем не менее и в ViT, и в SegFormer сложность внимания остаётся квадратичной относительно числа токенов $N$, 
 что затрудняет обработку изображений высокого разрешения и применение таких моделей в системах реального времени. 
 Актуальность работы обусловлена необходимостью совместить точность трансформерных сегментаторов с быстродействием, 
 достаточным для работы в реальном времени.

Научная проблема состоит в противоречии между высокой точностью сегментации, которую обеспечивает глобальный механизм внимания, 
и его квадратичной вычислительной сложностью, препятствующей работе в реальном времени.

Объектом исследования являются сегментаторы на основе визуальных трансформеров, предметом исследования - 
методы линеаризации механизма внимания в таких сегментаторах. Целью работы является поиск метода линеаризации внимания, 
который сохраняет точность сегментации ViT-сегментатора и позволяет ему работать в реальном времени. 
Для достижения цели решаются задачи анализа существующих методов линеаризации внимания, их сравнения по
 точности и быстродействию на задачах сегментации, выбора и адаптации наиболее подходящего метода к многоуровневой архитектуре сегментатора, 
 а также его экспериментальной апробации.

Результаты работы применимы в сегментации медицинских снимков, в том числе при выделении органов и новообразований 
на снимках компьютерной томографии, в анализе дорожных сцен для систем помощи водителю и беспилотного транспорта, в 
обработке спутниковых и аэрофотоснимков, в промышленном визуальном контроле дефектов, а также при построении 
динамической зоны интереса в области трамвайных путей \cite{myarticle}.

\subsection{Введение}
	Появление механизма внимания и трансформеров, изначально предложенных для задач обработки 
	естественного языка, привело к их последующей интеграции в задачи компьютерного зрения, в том числе для задач 
	семантической и экземплярной сегментации. Ключевым преимуществом таких моделей является способность
	моделировать глобальные и локальные зависимости между удаленными областями изображения уже
	на ранних стадиях обработки тогда как в сверточных архитектурах глобальный
	контекст формируется лишь с помощью последовательных операций сверток.

	Ограничение сверточных архитектур видно при попытке сегментировать мелкогабаритные или тонкие протяженные объекты, такие как
	ЖД и трамвайные рельсы, так как они занимают малую долю площади всего кадра и имеют изменчивую ширину и длину на кадрах \cite{myarticle}. Помимо этого, 
	их признаки существенно изменяются из-за искожений кадра, бликов, теней и окклюзий, в последствии чего, энкодер CNN приводит к утрате информации, из-за 
	чего маски сегмента могут получится либо оборванными либо не накладываться вообще. 
	Экземплярная сегментация так же не решает такую проблему. Например сети унаследывавшие принципт YOLACT, которые используют наложение итоговой маски как линейную 
	комбинацию прототипных карт с последующим обрезанием по ограничивающей рамки так же не справляются с моментами когда тень падает на рельсы или на кадрах есть блики. 
	Помимо этого ограничивающая рамка может обрезать кусок маски сегмента если ветка детекции не уловит рельсы целеком \cite{myarticle}. 

	С внедрением архитектур на основе Vision Transformer (ViT), точность сегментации для тонких и мелкогабаритных объектов заметно выросла, что обусловило 
	применение трансформерных сегментаторов и детекторов в задачах компьютерного зрения. 
	
	Модель SegFormer \cite{segformer}, построенная на механизме эффективного самовнимания
	(Efficient Self-Attention, ESA), была представлена наименьшей c конфигурацией
	B0, содержащей приблизительно 3.8 млн обучаемых параметров. Несмотря на столь
	малое число параметров, данная модель показала конкурентоспособные значения
	метрик качества. Это позволяет предположить, что переход к более емким конфигурациям семейства B3 и B5 обеспечил бы прирост mIoU и mDice, 
	заметно превосходящий показатели рассмотренных сверточных архитектур.

	Однако используемый в SegFormer механизм ESA имеет сложность $O(\frac{N^2}{R})$, где $N$ - число пикселей, а $R$ - коэффициент понижения размерности.

	Вся проблема заключается в квадратичной сложности, что дает сильное ограничение на скорость инференса и обучение модели.
	В литературе активно развивается направление линеаризации механизма внимания. Целью таких исследований является снижение асимптотической сложности до $O(N)$,
	при сохранении способности моделировать глобальные зависимости. Отсюда следует гипотеза исследования - линеаризация механизма внимания в архитектуре 
	SegFormer позволяющая получить сегментатор сочитающий быструю скорость инференса сравнимую для работы в реальном времени и сохраняющую точность свойственную трансформерной архитектре.


\subsection{Проблемы современных ViT}

\subsection{Постановка задачи}

\subsection{Методы}

\subsection{Эксперемент}

\subsection{Результаты}

\subsection{Выводы}

\renewcommand{\refname}{Список источников}
\begin{thebibliography}{9}
\addcontentsline{toc}{section}{Список источников}

\bibitem{myarticle} Салимли~А.~М., Востров~А.~В. Построение зоны интереса на основе сегментации // Неделя науки ИКНК. - 2026. - С.~386.

\end{thebibliography}


\end{document}